\documentclass[12pt]{article}

\usepackage[utf8]{inputenc}
\usepackage[T1]{fontenc}
\usepackage[brazil]{babel}
\usepackage{graphicx}
\usepackage{geometry}
\usepackage{xcolor}
\usepackage{indentfirst}
\usepackage{listings}
\usepackage{hyperref}

\geometry{a4paper, left=3cm, top=3cm, right=2cm, bottom=2cm}

\hypersetup{
    colorlinks=true,
    linkcolor=blue,
    urlcolor=blue,
    citecolor=blue
}

\lstset{
    basicstyle=\ttfamily\small,
    keywordstyle=\color{blue},
    stringstyle=\color{red},
    commentstyle=\color{green!60!black},
    breaklines=true,
    frame=single,
    captionpos=b
}

\newcommand{\placeholder}[1]{\textit{[#1]}}

\begin{document}

\begin{titlepage}
    \begin{center}
        \IfFileExists{logo_ufal.png}{%
            \includegraphics[width=0.2\textwidth]{logo_ufal.png}\\[0.5cm]
        }{%
            \vspace*{1.5cm}
        }
        {\large \textbf{UNIVERSIDADE FEDERAL DE ALAGOAS}} \\
        {\large \textbf{INSTITUTO DE COMPUTAÇÃO}} \\
        {\large \textbf{CIÊNCIA DA COMPUTAÇÃO}} \\
        \vfill
        {\Large \textbf{RELATÓRIO TÉCNICO FINAL INTEGRADO}} \\
        {\large \textbf{Framework para Compiladores e Interpretadores Usando Microsserviços}} \\
        \vfill
        \textbf{Equipe:}\\
        Bruno (Responsável pelas 3 disciplinas) \\
        Alan, Karlisson e Maria \\
        \vfill
        \textbf{Professor Orientador:}\\
        Dr.\ Arturo Hernandez Dominguez \\
        \vfill
        Maceió, AL \\
        2026
    \end{center}
\end{titlepage}

\tableofcontents
\newpage

\section{Introdução}
\subsection{Enunciado do Projeto}
Este relatório consolida o trabalho desenvolvido como requisito parcial integrado das disciplinas ministradas pelo Prof.\ Arturo Hernandez Dominguez no Instituto de Computação da UFAL, abrangendo Compiladores, Reuso de Software e Linhas de Produto de Software (LPS). O artefato central é o \emph{MiniPar Framework 2026.1}: uma linha de produto de software organizada em microsserviços, destinada ao processamento da linguagem MiniPar com orientação a objetos, construções de paralelismo \texttt{par}/\texttt{seq} e múltiplos back-ends (interpretador, C/C++, Rust e ARM).

A especificação acadêmica do framework \cite{minipar2026} deriva a evolução da MiniPar 2025.1 documentada por Maciel \cite{maciel2025} e pelos geradores procedimentais de Rego \cite{rego2025}. O repositório mantém documentação de gestão alinhada ao código: \texttt{COMPLIANCE\_AUDIT.md} (conformidade e gaps), \texttt{SCHEDULE.md} (cronograma), \texttt{ACTIVITIES.md} (sprints e checklist E2E) e \texttt{ROADMAP.md} (entregas por fase). Em vez de concentrar todas as fases de compilação em um único binário monolítico, adotou-se arquitetura distribuída em que cada etapa ou variante de saída corresponde a um serviço autônomo, comunicando-se por REST e JSON. Essa decisão articula, de um lado, a engenharia de domínio aplicada ao domínio ``compiladores MiniPar'' e, de outro, o desacoplamento típico de sistemas baseados em microsserviços.

\subsection{Objetivos}
Objetiva-se (i)~completar o suporte orientado a objetos na gramática e na AST da MiniPar 2026.1; (ii)~expor análise léxica, sintática e semântica como serviços HTTP consumíveis por um front-end e por um API Gateway; (iii)~materializar os back-ends da LPS (interpretador, geração para C/C++, Rust e ARM) como microsserviços que reutilizam um núcleo comum (\texttt{minipar-core}); e (iv)~registrar, na documentação e nos diagramas do repositório, os pontos de variação e as variantes que caracterizam a linha de produto.

\section{Metodologia Ágil e Planejamento}
\subsection{Scrum e Sprints}
O cronograma seguiu desenvolvimento incremental por sprints, com entrega prevista para 10 de junho de 2026. Na \textbf{Fase~1} (2--4/jun), priorizou-se o desbloqueio do pipeline real até a análise semântica: lexer, parser descendente recursivo, serialização da AST e verificação de classes e herança. Na \textbf{Fase~2} (3--5/jun), incorporou-se o padrão \emph{Template Method} no pacote \texttt{translation/}, os microsserviços de interpretação e de geração de código, e a orquestração do gateway com \texttt{PIPELINE\_BACKEND\_MODE=http}. Na \textbf{Fase~3} (5--7/jun), implementou-se o coordenador \texttt{ms-parallel-coord}, os \emph{workers} socket, o fractal de Sierpinski em MiniPar OO e a integração E2E do modo \texttt{DISTRIBUTED\_SOCKETS}.

\subsection{Backlogs}
\begin{itemize}
    \item \textbf{Product Backlog:} gestão da variabilidade LPS (modo interpretador \emph{versus} compilador, back-end alvo, ambiente local ou distribuído); casos de teste de paralelismo; demonstração do fractal.
    \item \textbf{Sprint Backlog Fase 1:} pacote \texttt{minipar-core}; microsserviços \texttt{ms-front-end} e \texttt{ms-semantic}; Docker Compose com \texttt{PIPELINE\_MODE=http}.
    \item \textbf{Sprint Backlog Fase 2:} módulo \texttt{translation/} (\emph{Template Method}); \texttt{ms-interpreter}; \texttt{ms-codegen-c}, \texttt{ms-codegen-rust}, \texttt{ms-codegen-arm}; variável \texttt{PIPELINE\_BACKEND\_MODE=http}.
    \item \textbf{Sprint Backlog Fase 3:} \texttt{ms-parallel-coord}; workers socket; fractal Sierpinski (\texttt{13\_sierpinski.minipar}); integração \texttt{DISTRIBUTED\_SOCKETS} E2E.
\end{itemize}
\subsection{Arquitetura de Microsserviços}
O ponto de entrada único é o \texttt{api-gateway} (porta~3000), que expõe \texttt{POST /api/v1/process} e encaminha o fluxo conforme a variabilidade informada pelo cliente. O serviço \texttt{ms-front-end} (:3001) realiza a análise léxica e sintática, devolvendo a AST em JSON; o \texttt{ms-semantic} (:3002) consome essa AST e produz a tabela de símbolos e eventuais diagnósticos semânticos. Os back-ends da LPS (:3003--3007) recebem a AST já validada e executam ou geram código conforme o valor de \texttt{targetVariability} selecionado na interface.

Do ponto de vista de engenharia de domínio, o domínio ``pipeline MiniPar'' foi particionado em subsistemas com contratos estáveis (\texttt{\_AST\_CONTRACT.md}), de modo que novas variantes de back-end possam ser acopladas sem alterar o lexer ou o analisador semântico. A Figura conceitual e os fluxos detalhados encontram-se em \texttt{docs/diagrams/}: \texttt{architecture.mmd}, \texttt{pipeline-sequence.mmd}, \texttt{template-method.mmd}, \texttt{codegen-c-flow.mmd}, \texttt{reuse-map.mmd} e \texttt{validation-cases.mmd}.

\subsection{Gestão da Variabilidade (LPS)}
Identificam-se três pontos de variação principais \cite{pohl2010}: (1)~modo de execução (interpretador ou compilador); (2)~back-end de código (C, C++, Rust, ARM); (3)~ambiente de execução (local ou distribuído via \texttt{ms-parallel-coord}). O diagrama de \emph{features} encontra-se em \texttt{docs/diagrams/feature-tree.mmd}. Cada combinação válida mapeia para um subconjunto de microsserviços invocado pelo gateway.

\section{Modelagem UML}
\subsection{Diagrama de Casos de Uso}
Os casos de uso principais são: editar código MiniPar, selecionar variante LPS, executar pipeline, visualizar AST e tabela de símbolos, executar paralelismo distribuído (3 máquinas) e visualizar o fractal de Sierpinski. O diagrama encontra-se em \texttt{docs/diagrams/uml-use-cases.mmd}.

\subsection{Arquitetura de Componentes}
Na modelagem por componentes, Léxico, Parser, Analisador Semântico e Geradores aparecem como unidades substituíveis encapsuladas em microsserviços FastAPI, orquestradas pelo \texttt{api-gateway} NestJS. O diagrama UML de componentes está em \texttt{docs/diagrams/uml-components.mmd}. Na Fase~1, concentrou-se a lógica de análise em \texttt{packages/minipar-core}. Na Fase~2, o subsistema de tradução (\texttt{minipar\_core/translation}) agrega TAC, interpretador e emissores de código, todos subordinados ao esqueleto \texttt{AbstractBackendTranslator}. Na Fase~3, adicionou-se o subsistema de paralelismo (\texttt{ms-parallel-coord} + workers).

\subsection{Diagrama de Classes --- framework e AST}
O diagrama de classes do framework (Template Method, Lexer, Parser, back-ends) encontra-se em \texttt{docs/diagrams/uml-classes-framework.mmd}. A AST orientada a objetos replica, em JSON, a estrutura dos nós Java de Maciel \cite{maciel2025}. Os nós centrais são:
\begin{itemize}
    \item \texttt{Program} --- raiz; atributo \texttt{declarations[]}
    \item \texttt{ClassDecl} --- \texttt{name}, \texttt{extends}, \texttt{members[]}
    \item \texttt{MethodDecl} --- \texttt{returnType}, \texttt{name}, \texttt{parameters}, \texttt{body}
    \item \texttt{ParBlock} / \texttt{SeqBlock} --- paralelismo estruturado e sequência
    \item \texttt{NewInstance}, \texttt{MethodCall}, \texttt{PrintStmt} --- instanciação, despacho de métodos e saída
\end{itemize}

\section{Fases do Compilador MiniPar 2026.1}
\subsection{Gramática BNF}
A gramática evolui a MiniPar 2025.1 \cite{maciel2025} com orientação a objetos e paralelismo. Produções principais:

\begin{lstlisting}[caption={BNF MiniPar 2026.1 OO (extrato)}]
<programa>      ::= { <declaracao> }
<declaracao>    ::= <classe> | <funcao> | <canal> | <stmt>
<classe>        ::= "class" ID [ "extends" ID ] "{" { <membro> } "}"
<membro>        ::= <atributo_oo> | <metodo>
<atributo_oo>   ::= <tipo> ID [ "=" <expr> ] ";"
<metodo>        ::= <tipo> ID "(" [ <param> { "," <param> } ] ")" <bloco>
<stmt>          ::= <bloco> | <if> | <while> | <for> | <print> | <return>
                  | "par" "{" { <stmt> } "}" | "seq" "{" { <stmt> } "}"
<expr>          ::= <atrib> | <logico_or>
<atrib>         ::= <logico_or> [ "=" <atrib> ]
<instanciacao>  ::= "new" ID "(" [ <expr> { "," <expr> } ] ")"
<chamada_met>   ::= <expr> "." ID "(" [ <expr> { "," <expr> } ] ")"
<tipo>          ::= "void" | "number" | "string" | "bool" | ID
\end{lstlisting}

Tokens de paralelismo: \texttt{par}, \texttt{seq}, \texttt{s\_channel}, \texttt{c\_channel}. Palavras reservadas OO: \texttt{class}, \texttt{extends}, \texttt{new}, \texttt{this}, \texttt{super}.
\subsection{Pseudocódigos dos Analisadores}
\subsubsection{Analisador Léxico}
O módulo \texttt{minipar\_core.lexer} implementa a classe \texttt{Lexer}, cujo método \texttt{tokenize()} percorre o código-fonte aplicando \texttt{skip\_whitespace()}, \texttt{skip\_comment()} e os leitores especializados abaixo.

\begin{lstlisting}[language=Python,caption={Fluxo de \texttt{Lexer.tokenize()} (minipar\_core.lexer)}]
enquanto pos < len(source):
    skip_whitespace()
    se skip_comment(): continuar
    se digito           -> read_number()      # NUMBER_LITERAL
    se aspas duplas     -> read_string()      # STRING_LITERAL
    se letra ou _       -> read_identifier()  # KEYWORD ou IDENTIFIER
    senao               -> operadores e delimitadores (#, ==, par, ...)
emitir Token(EOF)
\end{lstlisting}

\subsubsection{Analisador Sintático e AST}
O \texttt{Parser} descendente recursivo consome a lista de \texttt{Token} e constrói um \texttt{Program}. As produções de alto nível incluem \texttt{class\_declaration()}, \texttt{method\_declaration()}, \texttt{oo\_var\_declaration()} (tipo, nome e inicializador opcional), blocos \texttt{par}/\texttt{seq}, e comandos \texttt{print}/\texttt{println}. A serialização para intercâmbio entre microsserviços utiliza \texttt{ast\_json.to\_dict}.

\subsubsection{Coordenador de Paralelismo Distribuído}
\begin{lstlisting}[caption={Pseudocodigo ms-parallel-coord /coordinate}]
para cada host em [quicksort, matrix, factorial]:
    disparar thread dispatch_worker(host, port)
aguardar todas as threads
ordenar resultados por role
montar output agregado com rotulos PC1, PC2, PC3
retornar { output, results[] }
\end{lstlisting}

\subsubsection{Worker socket (processo independente)}
\begin{lstlisting}[caption={Pseudocodigo worker socket}]
servidor escuta na porta configurada (9001--9003)
ao receber conexao:
    ler comando "RUN"
    executar algoritmo (QuickSort | Matriz | Fatorial)
    enviar JSON { role, machine, data } via socket
    fechar conexao
\end{lstlisting}

\subsubsection{Analisador Semântico e Tabela de Símbolos}
A função \texttt{semantic\_json.analyze\_ast\_dict} percorre o dicionário da AST: registra declarações de classe, valida cláusulas \texttt{extends}, sinaliza membros duplicados e monta \texttt{symbolTable.scopes[]} para inspeção na UI e persistência opcional no PostgreSQL.

\section{Técnicas de Reuso}
\subsection{Padrão Template Method}
A variabilidade entre back-ends da LPS não foi tratada com ramificações dispersas no gateway, e sim com engenharia de domínio apoiada no padrão \emph{Template Method} \cite{gamma1995}. A classe \texttt{AbstractBackendTranslator}, em \texttt{packages/minipar-core/minipar\_core/translation/}, define o esqueleto fixo \texttt{translate(ast\_dict)}: \texttt{validate} $\rightarrow$ \texttt{prepare} $\rightarrow$ \texttt{emit} (hotspot) $\rightarrow$ \texttt{finalize}. Cada variante concreta especializa os hotspots de emissão e de materialização do resultado.

\begin{itemize}
    \item \texttt{InterpreterBackend} --- percorrência direta da AST, portando a semântica dinâmica de \texttt{cl-minipar} \cite{maciel2025};
    \item \texttt{CBackend} / \texttt{CppBackend} --- geração via TAC e invocação de \texttt{gcc -O2}, conforme a cadeia de Rego \cite{rego2025};
    \item \texttt{RustBackend} --- emissão mínima com \texttt{rustc};
    \item \texttt{ARMBackend} --- lowering TAC $\rightarrow$ ARMv7.
\end{itemize}

O diagrama de classes do padrão encontra-se em \texttt{docs/diagrams/template-method.mmd}; o mapa de artefatos reutilizados, em \texttt{reuse-map.mmd}.

\begin{lstlisting}[language=Python,caption={Esqueleto \texttt{AbstractBackendTranslator.translate} (base\_translator.py)}]
def translate(self, ast_dict):
    self.validate(ast_dict)
    if self._errors:
        return TranslationResult(..., errors=self._errors)
    self.prepare(ast_dict)
    self.emit(ast_dict)      # hotspot por back-end
    return self.finalize()
\end{lstlisting}

\subsection{Componentes Reutilizados}
O reuso sistemático partiu dos repositórios de referência da turma 2025.1, migrados para Python no núcleo compartilhado:

\begin{tabular}{|l|l|}
\hline
\textbf{Origem} & \textbf{Destino no framework} \\
\hline
\texttt{cl-minipar} (Lexer, Parser, AST OO, Interpreter) & \texttt{minipar-core} + \texttt{translation/interpreter.py} \\
\texttt{projeto\_compiladores} (TAC, C, ARM, gcc) & \texttt{translation/tac\_codegen.py}, \texttt{c\_backend.py}, \texttt{arm\_backend.py} \\
\texttt{projeto\_compiladores} (semântica procedural) & base de \texttt{semantic.py} / \texttt{semantic\_json.py} \\
\hline
\end{tabular}

\section{Implementação}
\subsection{Estrutura do Monorepo}
O repositório agrupa \texttt{frontend/} (Angular), \texttt{api-gateway/} (NestJS), \texttt{packages/minipar-core/}, \texttt{microservices/}, \texttt{database/} (scripts PostgreSQL) e \texttt{docs/diagrams/}.

\subsection{Tecnologias Utilizadas}
Front-end em Angular; gateway em NestJS; microsserviços de análise e código em FastAPI (Python~3.11); persistência de histórico de compilações em PostgreSQL; orquestração local via Docker Compose.

\subsection{API Gateway e Microsserviços}
\begin{tabular}{|l|l|l|l|}
\hline
\textbf{Serviço} & \textbf{Rota} & \textbf{Porta} & \textbf{Fase} \\
\hline
ms-front-end & POST /parse & 3001 & 1 \\
ms-semantic & POST /analyze & 3002 & 1 \\
ms-interpreter & POST /execute & 3003 & 2 \\
ms-codegen-c & POST /generate & 3004 & 2 \\
ms-codegen-rust & POST /generate & 3005 & 2 \\
ms-codegen-arm & POST /generate & 3007 & 2 \\
ms-parallel-coord & POST /coordinate & 3006 & 3 \\
worker-quicksort & socket & 9001 & 3 \\
worker-matrix & socket & 9002 & 3 \\
worker-factorial & socket & 9003 & 3 \\
api-gateway & POST /api/v1/process & 3000 & 1--3 \\
\hline
\end{tabular}

As variáveis de ambiente \texttt{PIPELINE\_MODE=http} e \texttt{PIPELINE\_BACKEND\_MODE=http}, definidas no Docker Compose da Fase~2, alternam o encaminhamento entre implementações embutidas e chamadas HTTP aos microsserviços.

\section{Testes e Resultados}
\subsection{Resultados da Fase 1}
Validou-se o pipeline de análise estática pela interface web (ação \textbf{Executar}, atalhos \texttt{Ctrl+Enter} e \texttt{F5}) e por requisições diretas ao gateway:

\begin{lstlisting}[caption={Exemplo curl --- AST e symbolTable reais}]
curl -X POST http://localhost:3000/api/v1/process \
  -H "Content-Type: application/json" \
  -d '{"sourceCode":"class Main { void run() { println(\"ok\"); } }",
       "targetVariability":"INTERPRETER","executionMode":"LOCAL"}'
\end{lstlisting}

A resposta contém \texttt{pipelineSteps} com \texttt{ms-front-end: parse} e \texttt{ms-semantic: analyze} (sem sufixo \texttt{mock}), além de nós \texttt{ClassDecl} na AST serializada.

\subsection{Casos de validação manual}
Os \emph{fixtures} em \texttt{sources/examples/} cobrem erros sintáticos, semânticos e casos válidos (ver \texttt{docs/diagrams/validation-cases.mmd}):

\begin{tabular}{|l|l|l|}
\hline
\textbf{Caso} & \textbf{Arquivo} & \textbf{Resultado} \\
\hline
Classe com herança válida & \texttt{02\_class\_extends} & \texttt{success: true} \\
\texttt{extends Arvore} inexistente & \texttt{04\_semantic\_extends\_unknown} & Erro semântico \\
\texttt{extends} sem nome & \texttt{05\_parse\_extends\_missing} & \texttt{Expected superclass name} \\
\texttt{exts} (palavra-chave inválida) & \texttt{06\_parse\_invalid\_keyword} & \texttt{Expected LBRACE} \\
Expressão solta \texttt{a} & \texttt{07\_expr\_only} & \texttt{ExprStmt} (análise OK) \\
\hline
\end{tabular}

\subsection{Casos de validação --- Fase 2 (back-ends)}
\begin{tabular}{|l|l|l|l|}
\hline
\textbf{Arquivo} & \textbf{LPS} & \textbf{Resultado} & \textbf{Pipeline step} \\
\hline
\texttt{08\_interpreter\_ok} & INTERPRETER & Saída \texttt{ok} & \texttt{ms-interpreter: execute} \\
\texttt{09\_oo\_new} & INTERPRETER & Saída \texttt{woof} & idem \\
\texttt{10\_par\_seq} & INTERPRETER & Saída \texttt{a b c} & idem \\
\texttt{11\_codegen\_c} & C & C + \texttt{gcc -O2} + stdout & \texttt{ms-codegen-c: generate} \\
\texttt{12\_codegen\_rust\_stub} & RUST & Código Rust MVP & \texttt{ms-codegen-rust: generate} \\
\hline
\end{tabular}

Mensagens representativas registradas na Fase~1:
\begin{itemize}
    \item \texttt{Semantic error: Superclass 'Arvore' not found for class 'Dog'}
    \item \texttt{Parser error at 5:20: Expected superclass name}
    \item \texttt{Parser error at 5:11: Expected LBRACE}
\end{itemize}

Erros de análise aparecem no painel \textbf{Console} da UI (chip de status e mensagem destacada); execuções bem-sucedidas habilitam as abas Saída, Símbolos e AST.

\subsection{Resultados da Fase 2 --- Back-ends reais}
Com \texttt{PIPELINE\_BACKEND\_MODE=http}, interpretador e geradores passaram a responder com saída produzida pelo núcleo \texttt{minipar-core}, e não por respostas simuladas:

\begin{lstlisting}[caption={Interpretador --- stdout real}]
# targetVariability: INTERPRETER
# output: "ok\n"
class Main { void run() { println("ok"); } }
\end{lstlisting}

\begin{lstlisting}[caption={Codegen C --- resposta E2E}]
curl -X POST http://localhost:3000/api/v1/process \
  -H "Content-Type: application/json" \
  -d '{"sourceCode":"class Main { void run() { println(\"ok\"); } }",
       "targetVariability":"C","executionMode":"LOCAL"}'
# pipelineSteps: ms-front-end: parse, ms-semantic: analyze,
#                ms-codegen-c: generate
# output: "Compiled with gcc -O2\nok\n"
# generatedCode: program.c gerado
\end{lstlisting}

Os arquivos \texttt{08\_interpreter\_ok} a \texttt{12\_codegen\_rust\_stub} documentam o comportamento observado. Para o back-end C, o fluxo completo (TAC, emissão, compilação) está em \texttt{codegen-c-flow.mmd}. As variantes Rust e ARM permanecem em nível MVP: geram código sintaticamente válido e retornam mensagem explícita quando \texttt{rustc} ou a \emph{toolchain} ARM não estão disponíveis no contêiner.

\subsection{Teste de Real Paralelismo (3 Computadores)}
Implementou-se o microsserviço \texttt{ms-parallel-coord} (:3006) e três \emph{workers} Python como processos independentes (\texttt{worker-quicksort}, \texttt{worker-matrix}, \texttt{worker-factorial}) comunicando-se exclusivamente via sockets TCP (:9001--9003). O modo \texttt{DISTRIBUTED\_SOCKETS} na UI dispara o coordenador após a análise semântica; os resultados agregados aparecem no Console e no campo \texttt{distributedResults} da resposta JSON.

\begin{lstlisting}[caption={curl --- teste 3 maquinas via gateway}]
curl -X POST http://localhost:3000/api/v1/process \
  -H "Content-Type: application/json" \
  -d '{"sourceCode":"class Main { void run() { println(\"ok\"); } }",
       "targetVariability":"INTERPRETER",
       "executionMode":"DISTRIBUTED_SOCKETS"}'
# pipelineSteps inclui ms-parallel-coord: coordinate
# output agrega QuickSort (PC1), Matriz (PC2), Fatorial (PC3)
\end{lstlisting}

O diagrama de sequência encontra-se em \texttt{docs/diagrams/sequence-3-machines.mmd}. Fixture: \texttt{sources/examples/14\_distributed\_menu.minipar}.

\subsection{Simulacao de Fractal}
Implementou-se o tapete de Sierpinski recursivo em MiniPar OO (\texttt{sources/examples/13\_sierpinski.minipar}): a classe \texttt{Sierpinski} utiliza o metodo recursivo \texttt{isBlack} e exibe uma matriz de caracteres \texttt{*} (solido) e \texttt{.} (vazio) no Console da UI via interpretador. Ordem~3 produz matriz $27 \times 27$.

\begin{lstlisting}[caption={Trecho do fractal MiniPar OO}]
class Sierpinski {
  number isBlack(number x, number y, number n) {
    if (n == 1) { return 1; }
    number m = n / 3;
    if (x >= m && x < 2 * m && y >= m && y < 2 * m) { return 0; }
    return this.isBlack(x % m, y % m, m);
  }
  void draw(number order) { /* imprime matriz linha a linha */ }
}
\end{lstlisting}

\subsection{Casos de validacao --- Fase 3}
\begin{tabular}{|l|l|l|}
\hline
\textbf{Arquivo} & \textbf{Modo} & \textbf{Resultado} \\
\hline
\texttt{13\_sierpinski} & INTERPRETER + LOCAL & Matriz fractal no Console \\
\texttt{14\_distributed\_menu} & INTERPRETER + DISTRIBUTED\_SOCKETS & Resultados PC1--PC3 \\
\texttt{09\_oo\_new} & C + LOCAL & \texttt{gcc -O2} + saída \texttt{woof} (OO codegen) \\
\hline
\end{tabular}

\section{Conclusao}
O \emph{MiniPar Framework 2026.1} integra as tres vertentes do projeto: pipeline de compilador MiniPar OO (lexer, parser, semantica, geradores), reuso sistematico via \emph{Template Method} e componentes reutilizados de 2025.1, e gestao de variabilidade LPS materializada em microsserviços REST orquestrados pelo \texttt{api-gateway}. As Fases~1--3 entregam pipeline E2E real, back-ends funcionais (interpretador, C com \texttt{gcc -O2} incluindo OO basico, Rust/ARM MVP), teste de paralelismo distribuido em tres workers socket e simulacao do tapete de Sierpinski na interface web.

\subsection{Link para o Repositorio (GitHub)}
\url{https://github.com/brunom/compiladores} (monorepo \texttt{minipar-framework}; ajustar URL final na entrega).

\subsection{Link para o Video de Apresentacao}
\url{https://minipar-framework.vercel.app} (UI de demonstracao; inserir URL do video backup na entrega).

\begin{thebibliography}{12}
\bibitem{maciel2025}
Maciel, E.\ \emph{MiniPar 2025.1} --- interpretador orientado a objetos. Implementação de referência \texttt{cl-minipar}, turma Compiladores, UFAL, 2025.

\bibitem{rego2025}
Rego, H.\ \emph{Projeto compiladores MiniPar} --- geradores de código C e ARM, representação intermediária TAC e integração com \texttt{gcc}. Repositório \texttt{projeto\_compiladores}, UFAL, 2025.

\bibitem{minipar2026}
Hernandez Dominguez, A.\ \emph{Especificação do projeto integrado MiniPar Framework 2026.1} --- framework de microsserviços, LPS e requisitos de demonstração (paralelismo distribuído e fractal). Documento da disciplina, UFAL, 2026.1.

\bibitem{pohl2010}
Pohl, K.; Böckle, G.; van der Linden, F.\ \emph{Software Product Line Engineering: Foundations, Principles, and Techniques}. Springer, 2010.

\bibitem{sommerville2016}
Sommerville, I.\ \emph{Software Engineering}. 10th ed. Pearson, 2016.

\bibitem{gamma1995}
Gamma, E. et al.\ \emph{Design Patterns: Elements of Reusable Object-Oriented Software}. Addison-Wesley, 1995.
\end{thebibliography}

\end{document}
